\section{Completed first probes}

Every hypothesis received a formal definition, an executable finite or
abstract probe, a hostile test, and a claim-level audit.  The complete output
is stored in \texttt{hypothesis\_probe\_results.json}.  Twelve standard-library
tests reproduce the calculations.

\begin{longtable}{p{0.07\textwidth}p{0.43\textwidth}p{0.36\textwidth}}
\toprule
ID & Probe result & Status\\
\midrule
\endhead
H1 & For \(Q=\left(\begin{smallmatrix}0&q^2\\1&0\end{smallmatrix}\right)\),
the positive weight \((q,1)\) gives weighted ceiling \(q=0.08\), while the
scalar row ceiling is one. & Finite positive-kernel mechanism survives; its
identification with a Yang--Mills vacuum weight is open.\\
\addlinespace
H2 & Uniform factor \(0.7\) gives rate \(-\log 0.7=0.356675\).  The strict
factors \(e^{-1/n}\) have zero asymptotic physical rate. & Pointwise strict
contraction is falsified as a sufficient premise.\\
\addlinespace
H3 & One star has no zero-spill modes.  The central four-star and eight-star
systems have respectively four and seven numerical zero-spill directions. &
Collective cancellation is plausible, but no boundary scaling or nonlinear
repair is proved.\\
\addlinespace
H4 & Two integer fluxes with identical center label and Pfaffian charge evolve
to different charges after the same added context. & Center plus total charge
is compositionally incomplete; a fuller integer carrier is necessary.\\
\addlinespace
H5 & The event identity \(-\log\mu(A)=D(\mu(\cdot|A)\Vert\mu)\) and data
processing hold in the finite model.  Charge five with multiplicity one
thousand yields bad probability \(0.8708\). & Free-energy formulation
survives; action-only suppression is falsified.\\
\addlinespace
H6 & Centering and orientation have commutator rank two and operator norm one.
The raw and oriented covariance spectra differ. & A real representation-path
defect is present; the curvature interpretation remains conjectural.\\
\addlinespace
H7 & Five thousand deterministic quaternion samples reconstruct with error
below \(1.4\times10^{-15}\) and radius below \(\pi/2\); \(-I\) is regular in
the negative-center chart. & The two-center \(\SU\) atlas supports
atlas-relative locality.\\
\addlinespace
H8 & A four-state regular/defect partition has lumpability defect one.  Two
refinement rounds isolate the four states and reduce the defect to zero. & A
finite completeness-detector algorithm is viable.\\
\addlinespace
H9 & Spectrum \(\{0,0.4,1.2\}\) gives true gap \(0.4\); an incomplete core
seeing only the high mode reports \(1.2\). & The spectral mechanism survives,
and density is shown indispensable.\\
\addlinespace
H10 & Coarse averaging and observable lifting recover the invariant
two-dimensional algebra exactly while erasing a gauge-copy antisymmetric
direction. & Observable-selective recovery exists in a toy model only.\\
\bottomrule
\end{longtable}

\subsection{The coordinated-spill signal}

The multi-star calculation uses the oriented centered star map \(A_P\) and
the external plaquette map \(B_P\).  For feasible star data \(y\), it
eliminates \(\ker A_P\) and computes
\begin{equation}
 \inf_{x:A_Px=y}\|B_Px\|^2.
\end{equation}
The resulting numerical spectra are summarized below.

\begin{center}
\begin{tabular}{lrrrr}
\toprule
Polymer & rank \(A_P\) & zero modes & minimum positive & maximum\\
\midrule
one star & 5 & 0 & 1.121212 & 1.635983\\
adjacent pair & 10 & 0 & 0.842105 & 2.181818\\
central four & 18 & 4 & 0.591234 & 2.255172\\
all eight interior links & 24 & 7 & 0.750000 & 1.428571\\
\bottomrule
\end{tabular}
\end{center}

The one-star row reproduces the exact orientation-correct spectrum.  The new
zero modes show that collective cancellation exists on this finite complex.
They do not show that arbitrary repair directions can be decomposed into such
modes, that nonzero modes have boundary rather than volume cost, or that a
finite nonlinear transformation has a controlled Jacobian.

\subsection{A compositional topology obstruction}

Let
\begin{align}
 n&=(-1,-1,-1,-1,-1,-1),\\
 n'&=(-1,-1,-1,-1,1,1).
\end{align}
Both have center label \((1,1,1,1,1,1)\) modulo two and Pfaffian charge one.
After adding the same context \(e_{01}=(1,0,0,0,0,0)\), their charges become
zero and two.  Hence the pair ``center label plus total Pfaffian charge'' is
not closed under contextual composition.  Retaining the full six-component
integer flux, or equivalent bilinear pairing data, is not optional in a
blocking calculus.
